\chapter{Multidimensional Data Modelling}
\begin{mybox}{Anggota Kelompok}
\begin{enumerate}
    \itemsep0em
    \item Citra Putri Saragih (25/574147/PPA/07239)
    \item Ade Dwi Putra (25/574144/PPA/07237)
\end{enumerate}
\end{mybox}

\section{Pendahuluan}
\subsection{Deskripsi Domain}
\subsection{Latar Belakang}
\subsection{Tujuan}

\section{Analisis Kebutuhan Bisnis}
\subsection{Stakeholder Bisnis}
\subsection{Pertanyaan Analitik Utama}
\subsection{Proses Bisnis Utama}
\subsection{Sumber Data}

\section{Perancangan Dimensional Modeling}
\subsection{Penentuan Grain}
\subsubsection{Grain Setiap Fact Table}
\subsubsection{Alasan Pemilihan Grain}
\subsubsection{Implikasi terhadap Granularitas dan Performa Query}
\subsection{Pemetaan Proses Bisnis}
\subsubsection{Business Process dan Fact Table}
\subsubsection{Conformed Dimensions}
\subsection{Perancangan Fact Table}
\subsubsection{Jenis Fact Table}
\subsubsection{Measures dan Klasifikasi Aditivitas}
\subsubsection{Foreign Keys}
\subsubsection{Degenerate Dimensions}
\subsection{Perancangan Dimension Table}
\subsubsection{Surrogate Keys dan Atribut Dimensi}
\subsubsection{Hierarki Dimensi}
\subsubsection{Role-Playing Dimensions}
\subsubsection{Junk Dimensions}
\subsubsection{Outrigger Dimensions}
\subsection{Slowly Changing Dimensions}
\subsubsection{Atribut yang Berubah}
\subsubsection{Strategi SCD dan Justifikasi}
\subsection{Diagram Skema Dimensional}
\subsubsection{Star Schema atau Snowflake Schema}
\subsubsection{Relasi dan Tipe Data Utama}

\section{Synthetic Data Generation}
\subsection{Strategi Pembuatan Data Sintetis}
\subsection{Skema Data Sumber}
\subsection{Implementasi Generator Data}
\subsection{Volume, Distribusi, dan Variasi Data}
\subsection{Outlier dan Skenario SCD}
\subsection{Contoh Data Mentah}

\section{Data Mart}
\subsection{Arsitektur dan Pipeline ETL/ELT}
\subsection{Skema Final Data Mart}
\subsection{Dokumentasi Tabel dan Kolom}
\subsection{Konsistensi Referensial}
\subsection{Koneksi ke BI Tools}
\subsection{Semantic Layer}

\section{Verifikasi dan Validasi}
\subsection{Query Analitik}
\subsection{Hasil Query}
\subsection{Data Quality Check}
\subsubsection{Pemeriksaan Duplikat}
\subsubsection{Pemeriksaan Nilai Null}
\subsubsection{Pemeriksaan Konsistensi Referensial}

\section{Kesimpulan dan Refleksi}
\subsection{Kesimpulan}
\subsection{Tantangan yang Dihadapi}
\subsection{Keputusan Desain Utama}
\subsection{Kesiapan Data Mart untuk Visualisasi}

\section{Referensi}
